% GENERATED FILE. Edit canonical lessons, then run npm run generate:books.
% canonical-source-hash: fnv1a64:1a2b2d877e3762c8
% canonical-lessons: SA-C549-sprhayati, SA-C549-lalasa, SA-C549-pratiksa, SA-C549-didrksa, SA-C549-samihate

\chapter{Longings}
\label{ch:sa-longings}
\hlchaptermodality{549}

\begin{chapteropening}
\textbf{By the end of this chapter:} \emph{I can talk about what I long for and what I would do.}

Complete the last lesson of chapter 549: I can talk about what I long for and what I would do.
\end{chapteropening}

\section[\texorpdfstring{\sk{स्पृहयति} (spṛhayati)}{स्पृहयति (spṛhayati)}]{\sk{स्पृहयति} (spṛhayati) — is eager for; longs for}

\label{lesson:SA-C549-sprhayati}



\begin{quote}
Before the new one: say the Sanskrit for flustered, then the Sanskrit for a refuge.
\end{quote}

\subsection*{You'll want to know: \sk{स्पृहयति}}
\textbf{\sk{स्पृहयति}} — \emph{spṛhayati} — \textquotedblleft{}is eager for; longs for\textquotedblright{}.

\begin{quote}
\sk{यदि} \sk{भवान्} \sk{ज्ञानाय} \sk{स्पृहयति}, \sk{तर्हि} \sk{प्रतिदिनं} \sk{पठतु।}
\end{quote}

\begin{quote}
\emph{yadi bhavān jñānāya spṛhayati, tarhi pratidinaṃ paṭhatu.}
\end{quote}

\begin{quote}
If you long for knowledge, read every day.
\end{quote}

Past (he …): \sk{सः} \sk{ज्ञानाय} \sk{स्पृहितवान्}, \emph{saḥ jñānāya spṛhitavān}.

\begin{cousinweb}
\textbf{\sk{स्पृहयति}} (\emph{spṛhayati}) comes from the root √spṛh, \textquotedblleft{}to long for, to desire\textquotedblright{}, which takes the dative; MW compares Greek \emph{sperchō} and English \emph{spring} (MW).
\end{cousinweb}

\begin{grammarlens}[title={What you've built}]
One more word for saying what happens if.
\end{grammarlens}

\subsection*{Your turn}
\begin{itemize}
\raggedright
  \item \emph{Say it:} \emph{spṛhayati}
  \item \emph{Say it:} \emph{spṛhayati}, once more
  \item \emph{Say it:} \emph{yadi bhavān jñānāya spṛhayati, tarhi pratidinaṃ paṭhatu.} — If you long for knowledge, read every day.
  \item \emph{Recall:} say the Sanskrit for flustered, then the Sanskrit for a refuge, then say \emph{spṛhayati} again
\end{itemize}

\begin{tcolorbox}[breakable,colback=teal!4,colframe=teal!35!black,title={Before you move on}]
What does \sk{स्पृहयति} mean? (\textquotedblleft{}Is eager for; longs for\textquotedblright{}.) Say it once more.
\end{tcolorbox}
\section[\texorpdfstring{\sk{लालसा} (lālasā)}{लालसा (lālasā)}]{\sk{लालसा} (lālasā) — a craving}

\label{lesson:SA-C549-lalasa}



\begin{quote}
Before the new one: say the Sanskrit for a refuge, then the Sanskrit for is eager for.
\end{quote}

\subsection*{You'll want to know: \sk{लालसा}}
\textbf{\sk{लालसा}} — \emph{lālasā} — \textquotedblleft{}a craving\textquotedblright{}.

\begin{quote}
\sk{यदि} \sk{मधुरे} \sk{लालसा} \sk{अस्ति}, \sk{तर्हि} \sk{एकं} \sk{मोदकं} \sk{खादतु।}
\end{quote}

\begin{quote}
\emph{yadi madhure lālasā asti, tarhi ekaṃ modakaṃ khādatu.}
\end{quote}

\begin{quote}
If you crave something sweet, eat a sweet.
\end{quote}

A feminine noun.

\begin{cousinweb}
\textbf{\sk{लालसा}} (\emph{lālasā}) comes from the intensive of the root √las, \textquotedblleft{}to play, to shine\textquotedblright{}, as MW derives it; MW gives \emph{lālasa} and \emph{lālasā} as longing, ardent desire (MW).
\end{cousinweb}

\begin{grammarlens}[title={What you've built}]
One more word for saying what happens if.
\end{grammarlens}

\subsection*{Your turn}
\begin{itemize}
\raggedright
  \item \emph{Say it:} \emph{lālasā}
  \item \emph{Say it:} \emph{lālasā}, once more
  \item \emph{Say it:} \emph{yadi madhure lālasā asti, tarhi ekaṃ modakaṃ khādatu.} — If you crave something sweet, eat a sweet.
  \item \emph{Recall:} say the Sanskrit for a refuge, then the Sanskrit for is eager for, then say \emph{lālasā} again
\end{itemize}

\begin{tcolorbox}[breakable,colback=teal!4,colframe=teal!35!black,title={Before you move on}]
What does \sk{लालसा} mean? (\textquotedblleft{}A craving\textquotedblright{}.) Say it once more.
\end{tcolorbox}
\section[\texorpdfstring{\sk{प्रतीक्षा} (pratīkṣā)}{प्रतीक्षा (pratīkṣā)}]{\sk{प्रतीक्षा} (pratīkṣā) — a wait; expectation}

\label{lesson:SA-C549-pratiksa}



\begin{quote}
Before the new one: say the Sanskrit for is eager for, then the Sanskrit for a craving.
\end{quote}

\subsection*{You'll want to know: \sk{प्रतीक्षा}}
\textbf{\sk{प्रतीक्षा}} — \emph{pratīkṣā} — \textquotedblleft{}a wait; expectation\textquotedblright{}.

\begin{quote}
\sk{यदि} \sk{प्रतीक्षा} \sk{कठिना} \sk{अस्ति}, \sk{तर्हि} \sk{पुस्तकं} \sk{पठतु।}
\end{quote}

\begin{quote}
\emph{yadi pratīkṣā kaṭhinā asti, tarhi pustakaṃ paṭhatu.}
\end{quote}

\begin{quote}
If the wait is hard, read a book.
\end{quote}

A feminine noun.

\begin{cousinweb}
\textbf{\sk{प्रतीक्षा}} (\emph{pratīkṣā}) is built on \emph{prati}, \textquotedblleft{}towards\textquotedblright{}, and the root √īkṣ, \textquotedblleft{}to see\textquotedblright{}, which also gives \emph{pratīkṣate}, \textquotedblleft{}waits\textquotedblright{}; MW gives \emph{pratīkṣā} as expectation (MW).
\end{cousinweb}

\begin{grammarlens}[title={What you've built}]
One more word for saying what happens if.
\end{grammarlens}

\subsection*{Your turn}
\begin{itemize}
\raggedright
  \item \emph{Say it:} \emph{pratīkṣā}
  \item \emph{Say it:} \emph{pratīkṣā}, once more
  \item \emph{Say it:} \emph{yadi pratīkṣā kaṭhinā asti, tarhi pustakaṃ paṭhatu.} — If the wait is hard, read a book.
  \item \emph{Recall:} say the Sanskrit for is eager for, then the Sanskrit for a craving, then say \emph{pratīkṣā} again
\end{itemize}

\begin{tcolorbox}[breakable,colback=teal!4,colframe=teal!35!black,title={Before you move on}]
What does \sk{प्रतीक्षा} mean? (\textquotedblleft{}A wait; expectation\textquotedblright{}.) Say it once more.
\end{tcolorbox}
\section[\texorpdfstring{\sk{दिदृक्षा} (didṛkṣā)}{दिदृक्षा (didṛkṣā)}]{\sk{दिदृक्षा} (didṛkṣā) — a wish to see}

\label{lesson:SA-C549-didrksa}



\begin{quote}
Before the new one: say the Sanskrit for a craving, then the Sanskrit for a wait.
\end{quote}

\subsection*{You'll want to know: \sk{दिदृक्षा}}
\textbf{\sk{दिदृक्षा}} — \emph{didṛkṣā} — \textquotedblleft{}a wish to see\textquotedblright{}.

\begin{quote}
\sk{यदि} \sk{नगरस्य} \sk{दिदृक्षा} \sk{अस्ति}, \sk{तर्हि} \sk{श्वः} \sk{तत्र} \sk{गच्छामः।}
\end{quote}

\begin{quote}
\emph{yadi nagarasya didṛkṣā asti, tarhi śvaḥ tatra gacchāmaḥ.}
\end{quote}

\begin{quote}
If you want to see the city, we will go there tomorrow.
\end{quote}

A feminine noun.

\begin{cousinweb}
\textbf{\sk{दिदृक्षा}} (\emph{didṛkṣā}) comes from the desiderative of the root √dṛś, \textquotedblleft{}to see\textquotedblright{}: MW gives \emph{didṛkṣā} as the desire of seeing (MW).
\end{cousinweb}

\begin{grammarlens}[title={What you've built}]
One more word for saying what happens if.
\end{grammarlens}

\subsection*{Your turn}
\begin{itemize}
\raggedright
  \item \emph{Say it:} \emph{didṛkṣā}
  \item \emph{Say it:} \emph{didṛkṣā}, once more
  \item \emph{Say it:} \emph{yadi nagarasya didṛkṣā asti, tarhi śvaḥ tatra gacchāmaḥ.} — If you want to see the city, we will go there tomorrow.
  \item \emph{Recall:} say the Sanskrit for a craving, then the Sanskrit for a wait, then say \emph{didṛkṣā} again
\end{itemize}

\begin{tcolorbox}[breakable,colback=teal!4,colframe=teal!35!black,title={Before you move on}]
What does \sk{दिदृक्षा} mean? (\textquotedblleft{}A wish to see\textquotedblright{}.) Say it once more.
\end{tcolorbox}
\section[\texorpdfstring{\sk{समीहते} (samīhate)}{समीहते (samīhate)}]{\sk{समीहते} (samīhate) — strives for}

\label{lesson:SA-C549-samihate}



\begin{quote}
Before the new one: say the Sanskrit for a wait, then the Sanskrit for a wish to see.
\end{quote}

\subsection*{You'll want to know: \sk{समीहते}}
\textbf{\sk{समीहते}} — \emph{samīhate} — \textquotedblleft{}strives for\textquotedblright{}.

\begin{quote}
\sk{यदि} \sk{सः} \sk{सुखं} \sk{समीहते}, \sk{तर्हि} \sk{सः} \sk{अन्यस्य} \sk{साहाय्यं} \sk{करोतु।}
\end{quote}

\begin{quote}
\emph{yadi saḥ sukhaṃ samīhate, tarhi saḥ anyasya sāhāyyaṃ karotu.}
\end{quote}

\begin{quote}
If he strives for happiness, let him help others.
\end{quote}

Past (he …): \sk{सः} \sk{सुखं} \sk{समीहितवान्}, \emph{saḥ sukhaṃ samīhitavān}.

\begin{cousinweb}
\textbf{\sk{समीहते}} (\emph{samīhate}) is built on \emph{sam} and the root √īh, \textquotedblleft{}to strive\textquotedblright{}; MW gives \emph{samīh} as to strive after, to wish for, to endeavour to gain (MW).
\end{cousinweb}

\begin{grammarlens}[title={What you've built}]
One more word for saying what happens if.
\end{grammarlens}

\subsection*{Your turn}
\begin{itemize}
\raggedright
  \item \emph{Say it:} \emph{samīhate}
  \item \emph{Say it:} \emph{samīhate}, once more
  \item \emph{Say it:} \emph{yadi saḥ sukhaṃ samīhate, tarhi saḥ anyasya sāhāyyaṃ karotu.} — If he strives for happiness, let him help others.
  \item \emph{Recall:} say the Sanskrit for a wait, then the Sanskrit for a wish to see, then say \emph{samīhate} again
\end{itemize}

\begin{tcolorbox}[breakable,colback=teal!4,colframe=teal!35!black,title={Before you move on}]
What does \sk{समीहते} mean? (\textquotedblleft{}Strives for\textquotedblright{}.) Say it once more.
\end{tcolorbox}
